%================================
% note-setup-leftbar.tex
% fenglielie@qq.com 2026-08-21
%================================

\usepackage{amsmath,amsthm,amssymb}
\usepackage{mathtools}
\usepackage{mathrsfs}
\usepackage{bm}
\usepackage{extarrows}
\usepackage[a4paper, margin=1cm, footskip=0.5cm]{geometry}
\usepackage{float}
\usepackage{indentfirst}
\usepackage{booktabs,multirow}
\usepackage[shortlabels,inline]{enumitem}
\usepackage{appendix}

%%================================
%% FIX GLOBALE PER LISTE DOPO I TEOREMI
%%================================
% Risolve il "primo punto a destra del titolo" quando itemize/enumerate 
% seguono immediatamente l'apertura di un ambiente theorem-like.
\makeatletter
\AddToHook{env/itemize/begin}{\if@inlabel\mbox{}\par\vspace{-\baselineskip}\fi}
\AddToHook{env/enumerate/begin}{\if@inlabel\mbox{}\par\vspace{-\baselineskip}\fi}
\makeatother
%%================================

\usepackage[dvipsnames]{xcolor}
\usepackage{pifont}
\usepackage{graphicx}
\graphicspath{
    {./figure/}{./figures/}{./image/}{./images/}{./graphic/}{./graphics/}{./picture/}{./pictures/}
}
\usepackage{subcaption}

\usepackage[ruled,linesnumbered,noline]{algorithm2e}
\usepackage{listings}
\lstdefinestyle{simpleStyle}{
    basicstyle=\ttfamily\small,
    breaklines=true,
    keywordstyle=\color{blue},
    identifierstyle=\color{black},
    stringstyle=\color{violet},
    commentstyle=\color[RGB]{34,139,34},
    showstringspaces=false,
    numbers=left,
    numbersep=2em,
    numberstyle=\footnotesize,
    frame=single,
    framesep=1em,
}
\lstset{style=simpleStyle}

\usepackage{hyperref}
\hypersetup{
    colorlinks=true,linkcolor=,urlcolor=cyan
}

\usepackage[most]{tcolorbox}
\usepackage{keytheorems}

\tcbset{
    theorem box/.style={
            enhanced,
            breakable,
            sharp corners,
            leftrule=2pt,
            rightrule=0pt,
            toprule=0pt,
            bottomrule=0pt,
        },
    theorem box rounded/.style={
            theorem box,
            rounded corners,
        }
}

%%================================
%% STILI UFFICIALI KEYTHEOREMS
%%================================
\newkeytheoremstyle{breakplain}{
    inherit-style=plain,
    headpunct={},
    break
}
\newkeytheoremstyle{breakdefinition}{
    inherit-style=definition,
    headpunct={},
    break
}
\newkeytheoremstyle{breakremark}{
    inherit-style=remark,
    headpunct={},
    break
}
\newkeytheoremstyle{notestyle}{
    headfont=\color{orange!80!black}\bfseries,
    bodyfont=\normalfont,
    spaceabove=3pt,
    spacebelow=3pt,
    headpunct={},
    break,
    qed=\ensuremath{\color{orange!80!black}\diamond}
}
\newkeytheoremstyle{solutionstyle}{
    headfont=\bfseries,
    bodyfont=\normalfont,
    spaceabove=3pt,
    spacebelow=3pt,
    headpunct={},
    break,
    qed=\ensuremath{\square}
}

%%================================
%% DEFINIZIONE AMBIENTI
%%================================

\newkeytheorem{theorem}[style=breakplain, name=Teorema, numbered=true, parent=section,
    tcolorbox-no-titlebar={theorem box, colframe=RoyalPurple, colback=RoyalPurple!8}]
\newkeytheorem{theorem*}[style=breakplain, name=Teorema, numbered=false,
    tcolorbox-no-titlebar={theorem box rounded, colframe=RoyalPurple, colback=RoyalPurple!8}]

\newkeytheorem{prop}[style=breakplain, name=Proposizione, numbered=true, sibling=theorem,
    tcolorbox-no-titlebar={theorem box, colframe=RoyalPurple, colback=RoyalPurple!8}]
\newkeytheorem{prop*}[style=breakplain, name=Proposizione, numbered=false,
    tcolorbox-no-titlebar={theorem box rounded, colframe=RoyalPurple, colback=RoyalPurple!8}]

\newkeytheorem{corollary}[style=breakplain, name=Corollario, numbered=true, sibling=theorem,
    tcolorbox-no-titlebar={theorem box, colframe=NavyBlue, colback=SkyBlue!8}]
\newkeytheorem{corollary*}[style=breakplain, name=Corollario, numbered=false,
    tcolorbox-no-titlebar={theorem box rounded, colframe=NavyBlue, colback=SkyBlue!8}]

\newkeytheorem{lemma}[style=breakplain, name=Lemma, numbered=true, sibling=theorem,
    tcolorbox-no-titlebar={theorem box, colframe=NavyBlue, colback=SkyBlue!8}]
\newkeytheorem{lemma*}[style=breakplain, name=Lemma, numbered=false,
    tcolorbox-no-titlebar={theorem box rounded, colframe=NavyBlue, colback=SkyBlue!8}]

\newkeytheorem{claim}[style=breakplain, name=Claim, numbered=true, sibling=theorem,
    tcolorbox-no-titlebar={theorem box, colframe=NavyBlue, colback=SkyBlue!8}]
\newkeytheorem{claim*}[style=breakplain, name=Claim, numbered=false,
    tcolorbox-no-titlebar={theorem box rounded, colframe=NavyBlue, colback=SkyBlue!8}]

\newkeytheorem{definition}[style=breakdefinition, name=Definizione, numbered=true, parent=section,
    tcolorbox-no-titlebar={theorem box, colframe=ForestGreen, colback=ForestGreen!5}]
\newkeytheorem{definition*}[style=breakdefinition, name=Definizione, numbered=false,
    tcolorbox-no-titlebar={theorem box rounded, colframe=ForestGreen, colback=ForestGreen!5}]

\newkeytheorem{example}[style=breakdefinition, name=Esempio, numbered=true, parent=section,
    tcolorbox-no-titlebar={theorem box, colframe=RawSienna, colback=RawSienna!5}]
\newkeytheorem{example*}[style=breakdefinition, name=Esempio, numbered=false,
    tcolorbox-no-titlebar={theorem box rounded, colframe=RawSienna, colback=RawSienna!5}]

\newkeytheorem{esercizio}[style=breakdefinition, name=Esercizio, numbered=true, parent=section,
    tcolorbox-no-titlebar={theorem box, colframe=WildStrawberry!30, colback=WildStrawberry!5}]
\newkeytheorem{esercizio*}[style=breakdefinition, name=Esercizio, numbered=false,
    tcolorbox-no-titlebar={theorem box rounded, colframe=WildStrawberry!30, colback=WildStrawberry!5}]

\newkeytheorem{oss}[style=breakremark, name=Osservazione, numbered=true, parent=section]
\newkeytheorem{oss*}[style=breakremark, name=Osservazione, numbered=false]

\newcommand{\notehandmark}{%
    \raisebox{0.05ex}{\large\ding{45}}\hspace{0.4em}%
}
\newkeytheorem{note}[style=notestyle, name={\notehandmark Nota}, numbered=true, parent=section]
\newkeytheorem{note*}[style=notestyle, name={\notehandmark Nota}, numbered=false]

\newkeytheorem{solution}[style=solutionstyle, name=Soluzione, numbered=true, parent=section]
\newkeytheorem{solution*}[style=solutionstyle, name=Soluzione, numbered=false]

%%================================
%% COVER (Invariata)
%%================================
\usepackage{titling}
\newcommand{\extrainfo}{}
\renewcommand{\extrainfo}[1]{\renewcommand{\extrainfocontent}{#1}}
\newcommand{\extrainfocontent}{}
\newcommand{\makecover}[1]{%
    \begingroup
    \hypersetup{pageanchor=false}%
    \begin{titlepage}
        \newgeometry{margin=0in}
        \parindent=0pt
        \includegraphics[width=\linewidth]{#1} % size = 1280*1024
        \vfill
        \begin{center}
            \parbox{0.618\textwidth}{%
                \raggedleft{\bfseries \Huge \thetitle} \\[0.6pt]
                \rule{0.618\textwidth}{4pt} \\
            }
        \end{center}
        \vfill
        \begin{center}
            \parbox{0.618\textwidth}{%
                \raggedleft\Large
                \begin{tabular}{r}
                    \theauthor \\
                    \thedate   \\
                \end{tabular}%
            }
        \end{center}
        \vfill
        \begin{center}
            \parbox[t]{0.7\textwidth}{\centering \itshape \extrainfocontent}
        \end{center}
        \vfill
    \end{titlepage}
    \restoregeometry
    \endgroup
}

%================================
% DIMOSTRAZIONI (amsthm nativo)
%================================

\renewcommand{\proofname}{Dimostrazione}

% Ridefinizione nativa e sicura di amsthm per forzare l'a capo e il titolo in grassetto nero.
\makeatletter
\renewenvironment{proof}[1][\proofname]{\par
  \pushQED{\qed}%
  \normalfont \topsep6\p@\@plus6\p@\relax
  \trivlist
  \item[\hskip\labelsep\bfseries #1]\ignorespaces\leavevmode\par\nobreak
}{%
  \popQED\endtrivlist\@endpefalse
}
\makeatother

% Wrappare l'ambiente proof in tcolorbox nel modo più sicuro previsto dal pacchetto stesso.
\tcolorboxenvironment{proof}{
    theorem box,
    colframe=TealBlue,
    colback=TealBlue!5
}

%================================
% PENCIL NOTE
%================================
% Versione inline pura che supporta andate a capo interne (doppio invio o \\)
\newcommand{\pencilnote}[1]{{\color{gray!80!black}#1}}